\documentclass[titlepage,a4paper]{jsarticle}
\usepackage{../../../sty/import}% 各種パッケージインポート
\usepackage{../../../sty/title}% タイトルページの変更

%% タイトルページの変数
% レポートタイトル
\title{心理学第三回授業出席表}
% 提出日
\expdate{\today}
% 科目名
\subject{心理学特論}
% 分野
\class{情報経営システム工学分野}
% 学年
\grade{M1}
% 学籍番号
\mynumber{24336488}
% 記述者
\author{本間三暉}

\begin{document}
% titleページ作成
\maketitle
\section*{keyword\#1}
合理的経済人
\section*{keyword\#2}
欲求階層説
\section{授業スライド17(「やる気」を削いでゆく要因を振り返る)をもとに，自分の「やる気」が削がれる要因や，モチベーションが著しく下がってしまう傾向があればまとめてみてください．}
自分のやる気が削がれる要因を整理すると，まず努力に対して結果が伴わないときに，強く意欲が低下する傾向がある．

また，何をすればよいのか分からず，それを知る方法も分からないまま続けるときにも，無力感が生じてモチベーションが下がりやすい．
さらに，疲労がたまっているときや十分に食事を取れていないときは，身体的な余裕のなさがそのまま気力の低下につながる．
加えて，必須ではないことのために動かなければならないときや，自分がその面白さを理解できないときにも，行動を起こす気持ちが弱くなる．

そのほか，一人だけで取り組んでいるときや，努力を否定されたときにも，やる気は大きく削がれる．
このように，自分は成果の見えにくさ，不明瞭さ，心身の余裕のなさ，周囲との関わりの薄さによってモチベーションが下がりやすいと考えられる．

\section{自分のモチベーションが下がってしまった局面（エピソード）と，なぜそのように下がってしまったのか？（理由，もしくは分析）を5W1Hを踏まえ，なるべく具体的にわかりやすく簡潔にまとめてみてください．}
% what, who, when, where, why, how
自分のモチベーションが下がってしまった場面の一つは，今年の春休みに友人数人と自宅でVCをつなぎながらゲームをしていたときである．
そのゲームにはランクシステムがあり，ランクを上げるためには多くの知識を集めたり，練習を重ねたりする必要があった．
しかし，いくら努力してもなかなかランクが上がらず，自分よりも努力していない人がランクを上げていたことに加え，ゲーム内で正しい指示をしたにもかかわらず，
ランクを理由にその指示を無視され，馬鹿にされたことで，モチベーションが大きく下がってしまった．

さらに，そのゲームは1試合あたりに多くの時間がかかるため，まとまった時間を割いていたにもかかわらず，リスペクトのない言葉を浴びたことで続ける意欲を失った．
その後，学校が始まってまとまった時間を取りづらくなったこともあり，それ以来そのゲームはアンインストールして触っていない．
% 参考文献
\begin{thebibliography}{99}
\bibitem{} 心理学特論 授業スライド
\end{thebibliography}

\end{document}